\begin{ouverture}[geometrie]{14}{Thalès : triangles emboîtés}{Partie 4 · Géométrie et espace}
{Alignement · Parallélisme · Rapports · Réciproque}
\objectif{Utiliser Thalès dans des triangles emboîtés.}
\objectif{Démontrer un parallélisme avec la réciproque.}
\objectif{Vérifier l’ordre des points et les côtés correspondants.}
\prerequis{Fractions, proportionnalité et équations simples.}
\end{ouverture}
\pagestyle{otMain}\essentiel\fontsize{11}{15}\selectfont
\phantomsection\label{cours-14-1}
\essSection{1}{Des côtés proportionnels}
\begin{propriete}[Théorème de Thalès]Dans le triangle ABC, M appartient à $[AB]$ et N à $[AC]$, avec M et N distincts de A. Si $(MN)$ est parallèle à $(BC)$, alors $\frac{AM}{AB}=\frac{AN}{AC}=\frac{MN}{BC}$. Les petits côtés sont tous au numérateur, les grands tous au dénominateur.\end{propriete}
\label{fig-cours-14-1-cours-thales}\figureGeo{cours-thales}

\begin{exemple}Si $AB=8$, $AC=10$, $AM=4$ cm et $(MN)\parallel(BC)$, alors $AN/10=4/8=1/2$, donc $AN=5$ cm. Il faut annoncer les alignements et le parallélisme avant les rapports.\end{exemple}
\begin{attention}$AM$ est une distance depuis A, pas le morceau MB. Si $AB=8$ et $MB=3$, alors $AM=8-3=5$. Un schéma avec des points nommés évite de mélanger ces segments.\end{attention}
\clearpage
\phantomsection\label{cours-14-2}
\essSection{2}{Prouver un parallélisme}
\begin{propriete}[Réciproque]Dans un triangle ABC, M est sur $[AB]$ et N sur $[AC]$, distincts de A. Si $AM/AB=AN/AC$, alors les droites $(MN)$ et $(BC)$ sont parallèles. Les alignements et l’ordre des points font partie des conditions.\end{propriete}
\begin{exemple}Si $AM=3$, $AB=9$, $AN=4$, $AC=12$ cm, les deux rapports valent $1/3$. Avec les appartenances aux segments, la réciproque donne $(MN)\parallel(BC)$. Si les deux rapports sont différents, ce parallélisme est impossible : il imposerait l’égalité par le théorème direct.\end{exemple}
\begin{savoirfaire}[Choisir le sens]Pour calculer une longueur, on utilise généralement le théorème direct à partir d’un parallélisme connu. Pour prouver le parallélisme, on utilise la réciproque à partir des rapports de longueurs. Écrire les deux rapports avec le même sommet de départ et comparer exactement.\end{savoirfaire}
\begin{attention}Cette édition étudie la configuration des triangles emboîtés attendue en Quatrième. La configuration en papillon et les triangles semblables seront repris en Troisième. Une simple impression de parallélisme sur le dessin n’est pas une hypothèse.\end{attention}

\clearpage\rubrique{Découverte}
\begin{blocLivre}{ex-14-1}\begin{exercice}[1]{Un rapport simple}
Dans ABC, $M\in[AB]$, $N\in[AC]$ et $(MN)\parallel(BC)$. $AB=8$, $AC=10$ et $AM=4$ cm. Calculer AN.
\label{fig-ex-14-1}\figureGeo{th-01}
\end{exercice}\end{blocLivre}
\begin{blocLivre}{ex-14-2}\begin{exercice}[1]{Le segment parallèle}
M et N sont sur $[AB]$ et $[AC]$, avec $(MN)\parallel(BC)$. $AB=12$, $AM=8$ et $BC=9$ cm. Calculer MN.
\label{fig-ex-14-2}\figureGeo{th-02}
\end{exercice}\end{blocLivre}
\begin{blocLivre}{ex-14-3}\begin{exercice}[1]{Retrouver le grand côté}
M et N sont sur les segments du triangle ABC et $(MN)\parallel(BC)$. $AM=3$, $AB=9$ et $MN=4$ cm. Calculer BC.
\label{fig-ex-14-3}\figureGeo{th-03}
\end{exercice}\end{blocLivre}
\begin{blocLivre}{ex-14-4}\begin{exercice}[1]{Une réciproque}
$M\in[AB]$, $N\in[AC]$, $AB=10$, $AM=6$, $AC=8$ et $AN=4{,}8$ cm. Démontrer $(MN)\parallel(BC)$.
\label{fig-ex-14-4}\figureGeo{th-04}
\end{exercice}\end{blocLivre}
\begin{blocLivre}{ex-14-5}\begin{exercice}[1]{Des rapports à simplifier}
$M\in[AB]$, $N\in[AC]$, $AB=9$, $AM=6$, $AC=12$ et $AN=8$ cm. Les droites MN et BC sont-elles parallèles ?
\label{fig-ex-14-5}\figureGeo{th-05}
\end{exercice}\end{blocLivre}

\rubrique{Entraînement}
\begin{blocLivre}{ex-14-6}\begin{exercice}[2]{Un parallélisme impossible}
$M\in[AB]$, $N\in[AC]$, $AB=10$, $AM=6$, $AC=15$ et $AN=8$ cm. Les droites MN et BC peuvent-elles être parallèles ?
\label{fig-ex-14-6}\figureGeo{th-06}
\end{exercice}\end{blocLivre}
\begin{blocLivre}{ex-14-7}\begin{exercice}[2]{Calculer le bon segment}
Dans ABC, $M\in[AB]$, $N\in[AC]$ et $(MN)\parallel(BC)$. $AB=12$, $BM=3$ et $AC=8$ cm. Calculer AN.
\label{fig-ex-14-7}\figureGeo{th-07}
\end{exercice}\end{blocLivre}
\begin{blocLivre}{ex-14-8}\begin{exercice}[2]{Un coefficient décimal}
Dans la configuration emboîtée, $(MN)\parallel(BC)$, $AB=10$, $AM=4$, $AC=20$ et $BC=15$ cm. Calculer AN et MN.
\label{fig-ex-14-8}\figureGeo{th-08}
\end{exercice}\end{blocLivre}
\begin{blocLivre}{ex-14-9}\begin{exercice}[2]{Un protocole de construction}
Tracer un triangle ABC avec $AB=9$ cm, $AC=12$ cm et $\widehat{BAC}=50^\circ$. Placer M à 3 cm de A sur AB, puis construire N pour que $(MN)\parallel(BC)$. Donner AN.
\label{fig-ex-14-9}\figureGeo{th-09}
\end{exercice}\end{blocLivre}
\begin{blocLivre}{ex-14-10}\begin{exercice}[2]{Une proportion inverse}
Dans les triangles emboîtés AMN et ABC avec $(MN)\parallel(BC)$, $AB=7$, $BC=14$ et $MN=5$ cm. Calculer AM.
\label{fig-ex-14-10}\figureGeo{th-10}
\end{exercice}\end{blocLivre}

\rubrique{Problèmes}
\begin{blocLivre}{ex-14-11}\begin{exercice}[3]{Une hauteur par proportion}
Dans un modèle d’ombres, A, M, B sont alignés au sol avec $AM=2$ m, $AB=10$ m. MN et BC sont verticaux, $MN=1{,}5$ m, et A, N, C sont alignés. Calculer BC.
\label{fig-ex-14-11}\figureGeo{th-11}
\end{exercice}\end{blocLivre}
\begin{blocLivre}{ex-14-12}\begin{exercice}[3]{Un morceau restant}
$AB=8$, $AM=3$, $AC=10$ cm, avec M et N sur les segments et $(MN)\parallel(BC)$. Calculer AN puis NC.
\label{fig-ex-14-12}\figureGeo{th-12}
\end{exercice}\end{blocLivre}
\begin{blocLivre}{ex-14-13}\begin{exercice}[3]{Justifier puis calculer}
$M\in[AB]$, $N\in[AC]$, $AB=10$, $AC=15$, $AM=4$, $AN=6$ et $BC=12$ cm. Démontrer le parallélisme puis calculer MN.
\label{fig-ex-14-13}\figureGeo{th-13}
\end{exercice}\end{blocLivre}
\begin{blocLivre}{ex-14-14}\begin{exercice}[3]{Partager un segment}
Sur une demi-droite issue de A, on place cinq intervalles égaux jusqu’à C. M est au troisième intervalle. La parallèle à CB passant par M coupe AB en N. Si $AB=10$ cm, calculer AN et expliquer l’intérêt de la construction.
\label{fig-ex-14-14}\figureGeo{th-14}
\end{exercice}\end{blocLivre}
\begin{blocLivre}{ex-14-15}\begin{exercice}[3]{Deux réductions emboîtées}
ABC est rectangle en A, avec $AB=12$ et $AC=16$ cm. Sur AB, $AM=3$ et $AP=9$ cm. Les parallèles à BC par M et P coupent AC en N et Q. Calculer AN, AQ et l’aire du quadrilatère MNQP, délimité dans l’ordre M, N, Q, P.
\label{fig-ex-14-15}\figureGeo{th-15}
\end{exercice}\end{blocLivre}

\clearpage
